\documentclass[12pt,a4paper]{article}
\usepackage[UTF8]{ctex}
\usepackage{amsmath}
\usepackage{enumitem}
\usepackage{geometry}
\geometry{left=2.5cm, right=2.5cm, top=2.5cm, bottom=2.5cm}

\title{\textbf{问题假设五句总结（摘要/模型节素材）}}
\author{2026 年第七届"华数杯"大学生数学建模竞赛 B 题}
\date{2026/08/10}

\begin{document}
\maketitle

\section*{问题假设}

\begin{enumerate}[label=\arabic*., leftmargin=2em]
    \item 所有模块均为硬模块，尺寸固定、不可缩放变形，仅允许整体
    旋转（问题一至三为 $90^\circ$，问题四为四向），任意两模块互不
    重叠但允许紧贴放置，坐标取整数网格。
    \item 不考虑布线拥塞、热分布、电源网格等物理约束，竞赛提供的
    模块、线网与端子数据完整可靠，直接采用。
    \item 问题一在自由轮廓下以面积最小为主目标、面积相同时长宽比
    越接近 1 越优；问题二在固定轮廓
    $\text{side}=\lceil\sqrt{\Sigma A(1+0.15)}\rceil$ 内以半周长线长
    度量线长代价，引脚取几何中心近似、端子位置固定。
    \item 问题三中可行性判定随死区比例单调，并行多种子判定抑制
    启发式假阴性，$d^*$ 限定为"当前搜索强度下验证过的最小可行值"。
    \item 问题四中 L/T 型模块可精确表示为固定相对位置的矩形并集，
    无重叠判定在分解矩形层面进行，穷举完备性与下界可达性仅对本
    题 4 模块实例成立。
\end{enumerate}

\end{document}
